\clearpage
\onecolumngrid
\pdfbookmark[1]{Personal Philosophical Opinion}{personal-philosophical-opinion}
\section*{Personal Philosophical Opinion}

Regardless of how extensive and precise the data collected by our successors in the future may be, the equations with which they attempt to understand the universe will most likely remain primarily research models for estimating the expected number of worlds. It is difficult to expect them to provide an exact inventory of the actual state of affairs. With them, we will be able to connect data, test our assumptions, and determine the range within which the actual number of worlds most likely lies. Yet because of the limitations of observation, we will probably never arrive at a completely precise answer. Even with substantially more advanced technology, it would be almost impossible to examine every star in the Galaxy and obtain complete data for each one about its planets, their atmospheres, geological histories, and possible biospheres. For stars in other galaxies, such a task would be incomparably more demanding.

The scale of the problem can already be illustrated by a simple thought experiment. If we assume that the Milky Way contains at least 100 billion stars, then our successors, if they wished to physically examine every one of them within a thousand years, would have to investigate approximately 274,000 stars per day on average. Even the hypothetical possibility of travelling at or beyond the speed of light would solve only part of the problem. They would still need an enormous number of research devices, would have to coordinate their trajectories, provide energy, and devote enough time to observing each planetary system. A complete empirical survey of the Galaxy would therefore be extraordinarily difficult, and when it comes to the entire universe, such an undertaking is even harder to imagine.

Part of the answer to the apparent absence of other civilizations may already lie in these scales. If a planet with conditions suitable for the long-term preservation of life occurred around only one in several thousand stars, the number of interesting targets would quickly decrease. Astronomical habitability is, moreover, only the beginning of a long and poorly understood chain of events. We do not know the probability of the origin of life, the development of complex organisms, the emergence of intelligence, or the development of a technological civilization, nor do we know what proportion of such civilizations survive long enough to become interstellar. There could also be an additional civilizational selection filter, connected with the ability to limit violence, internal conflicts, and self-destructive potential, which increases together with technological power. Nor do we know for how long during its existence such a civilization remains detectable from afar.

Distances represent only one side of the problem. The other is time. Two civilizations may arise in the same Galaxy, yet be separated by a million years. One may disappear before the other even emerges. They never meet, and it is entirely possible that they never detect one another. The apparent silence of the universe therefore does not yet allow us to conclude that the universe is empty. The reason may lie in the rarity of life, immense distances, the short period during which technological civilizations are detectable, or simply in the limitations of our observational methods.

All of this raises an even deeper question. What would it actually mean for humanity to learn that millions of technological civilizations currently exist in the universe, perhaps including some that are more intelligent than we are? Such a discovery would undoubtedly represent one of the greatest epistemological, anthropological, and existential turning points in the history of humanity. It is difficult, however, to know what such knowledge would do to the human being. We would answer one of the oldest questions concerning our own existence and satisfy part of our desire for knowledge, while at the same time being forced to reconsider our conception of our own position in the universe. Humanity would become one among potentially millions of independent manifestations of consciousness and intelligence.

Human life itself would not thereby lose its value. What would change above all would be our conception of our own exceptional status. If we discovered that intelligence exists on millions of other worlds, we would be less rare as a phenomenon, but our lives would not thereby become any less real or important. The value of humanity is, after all, not determined solely by its rarity in the universe. It is also determined by our capacities for understanding, creating, preserving, and giving to others, as well as by the way in which we connect our knowledge with moral and ethical responsibility.

Perhaps, therefore, the most important result of such equations will not be a definitive answer to the question of how many hypothetical ExoEarth candidates exist, but rather a clearer recognition of the limits of our knowledge and of the responsibility we bear as inhabitants of what is currently the only planet with an active biosphere that we know with certainty exists. Behind life on Earth stand billions of years of successful evolutionary history, while its most recent result is a technological civilization whose own future, however, remains subject to very great uncertainty\ldots
